\documentclass[10pt,letterpaper,final]{report}
\usepackage[margin=.75in]{geometry}
\usepackage[utf8]{inputenc}
\usepackage{amsmath}
\usepackage{amsfonts}
\usepackage{amssymb}
\usepackage{amsthm}
\renewcommand\qedsymbol{$\blacksquare$}
\usepackage{enumerate}
\usepackage{enumitem}
\usepackage{hyperref}
\usepackage{pdfpages}
\usepackage{graphics}
\usepackage{graphicx}
\usepackage{tikz}
\usepackage{tikz-qtree}
\usepackage{forest}
\usepackage{float}
\usetikzlibrary{automata,arrows}
\usetikzlibrary{shapes.multipart, positioning}
\usepackage{mdframed}
\usepackage[
  backend=biber,
  style=numeric-comp,
  sorting=none
]{biblatex}
\addbibresource{references.bib}

\usepackage{minted}
\usemintedstyle{algol_nu}
\usepackage{xltabular}
\usepackage{pdflscape}
\usepackage{tcolorbox}


% Based on template from COSC-417 at Towson University, originally authored by Marius Zimand

\begin{document}
\begin{center}
  \fbox{
    \vbox{
      \begin{flushleft}
        Jonathan Llovet \\  % authors' names
        EN.605.202.81 \\  %class
        2026-07-21 \\  % date
        Module 7 - Homework 7 \\ % ID
      \end{flushleft}
      \center{\Large{\textbf{Data Structures - More on Trees and Sorting}}}
      %\end{mdframed}
    } % end vbox
  } % end fbox
  \vline
\end{center}

\medskip

\noindent\textbf{Problem 1:}
\medskip

Implement maketree, setleft, and setright for right in-threaded binary trees using the sequential array representation.

\medskip
\noindent\textbf{Answer:}
\medskip

\begin{minted}[linenos,breaklines,breakanywhere=true,fontsize={\fontsize{8pt}{9pt}\selectfont}]{Go}
struct Node {
    # no pointers needed because we can calculate them relative to a node's index
    data: T
    right_thread: int # index of the right thread successor
    right_is_thread: bool
}

function extend_array(array) {
    array = [] if array == null else array
    size = 1 if len(array) == 0 else len(array)
    new_array = []
    new_array.extend([null]*size*2)
    for i in range(size) {
        new_array[i] = array[i]
    }
    return new_array
}

function maketree(item: Any, tree_array: List, index: int) -> List {
    tree_array = [null] if tree_array == null else tree_array
    index = 0 if index == null else index
    node = Node(data=item)
    node.data = item
    node.right_thread = null
    if index >= len(tree_array) {
        tree_array = extend_array(tree_array)
    }
    tree_array[index] = node
    return tree_array
}


function parent_index(index: int) -> int | null {
    if index == 0 {
        return null
    }
    reduction = 2 if index % 2 == 0 else 1
    parent = (index - reduction) // 2
    return parent
}


function left_child_index(tree_array: List, index: int, write: bool = False) -> Tuple[List, int | null] {
    if index >= len(tree_array) {   
        return tree_array, null
    }
    left_child_index = (index * 2) + 1
    if left_child_index >= len(tree_array) {
        if write {
            tree_array = extend_array(tree_array)
            return tree_array, left_child_index
        }
        else {
            return tree_array, null
        }
    }
    elif left_child_index < len(tree_array) {
        return tree_array, left_child_index
    }
    else { # something went wrong
        return tree_array, null
    }
}


function right_child_index(tree_array: List, index: int, write: bool = False) -> Tuple[List, int | null] {
    if index >= len(tree_array) {
        return tree_array, null
    }
    right_child_index = (index * 2) + 2
    if right_child_index >= len(tree_array) {
        if write {
            tree_array = extend_array(tree_array)
            return tree_array, right_child_index
        }
        else {
            return tree_array, null
        }
    }
    elif right_child_index < len(tree_array) {
        return tree_array, right_child_index
    }
    else { # something went wrong
        return tree_array, null
    }
}


function setleft(item: Any, tree_array: List, index: int) -> List {
    if index >= len(tree_array) {
        error("Couldn't find node: That position doesn't exist in the tree")
    }
    if tree_array[index] == null {
        error("Couldn't find node: That node doesn't exist in the tree")
    }
    if not isinstance(tree_array[index], Node){
        error(f"The tree was corrupted. Can't proceed. {tree_array[index]} is not a Node.")
    }
    
    tree_array, lc_index = left_child_index(tree_array, index, write=True)
    if lc_index == null { # left child index is out of range
        error("Couldn't find node: the provided position doesn't exist in the tree")
    }
    # index is within range, possible that we have a node there
    if tree_array[lc_index] != null { # found a node alredy exists
        error("Cant set left child. The selected node already has a left child.")
    }
    if tree_array[lc_index] == null {
        lc_node = Node(item)
        # set the current node's (i.e. the left child's parent) as the right thread
        lc_node.right_thread = index
        lc_node.right_is_thread = True
        tree_array[lc_index] = lc_node
    }
    return tree_array
}


function setright(item: Any, tree_array: List, index: int) -> List {
    if index >= len(tree_array){
        error("Couldn't find node: That position doesn't exist in the tree")
    }
    if tree_array[index] == null{
        error("Couldn't find node: That node doesn't exist in the tree")
    }
    if not isinstance(tree_array[index], Node){
        error(f"The tree was corrupted. Can't proceed. {tree_array[index]} is not a Node.")
    }
    parent_node = tree_array[index]
    tree_array, rc_index = right_child_index(tree_array, index, write=True)
    if rc_index == null { # right child index is out of range
        error("Couldn't find node: the provided position doesn't exist in the tree")
    }
    
    # index is within range, possible that we have a node there
    if tree_array[rc_index] != null{ # found a node alredy exists
        error("Cant set right child. The selected node already has a right child.")
    }
    if tree_array[rc_index] == null{
        rc_node = Node(item)
        # copy the right thread pointer from the parent
        rc_node.right_thread = parent_node.right_thread
        rc_node.right_is_thread = True
        # parent's successor thread points to the right child
        parent_node.right_thread = rc_index
        parent_node.right_is_thread = False
        tree_array[index] = parent_node
        tree_array[rc_index] = rc_node
    }
    return tree_array
}
\end{minted}

\medskip
\noindent\textbf{Problem 2:}
\medskip

Implement inorder traversal for the right in-thread tree in the previous problem.

\medskip
\noindent\textbf{Answer:}
\medskip

\begin{minted}[linenos,breaklines,breakanywhere=true,fontsize={\fontsize{8pt}{9pt}\selectfont}]{Go}
function get_leftmost_descendant(tree_array, index) -> Tuple[Node, int] {
    tree_array, lc_index = left_child_index(tree_array, index, write=False)
    if lc_index == null {
        return tree_array[index], index
    }
    # left subtree
    leftmost_descendant = null
    while True {
        if lc_index == null {
            return leftmost_descendant, index
        }
        if tree_array[lc_index] == null {
            return leftmost_descendant, index
        }
        index = lc_index
        leftmost_descendant = tree_array[index]
        tree_array, lc_index = left_child_index(tree_array, index, write=False) 
    }
}


function traverse_inorder(tree_array: List, index: int = 0, visit=print) {
    if len(tree_array) == 0 {
        return
    }
    if len(tree_array) == 1 {
        visit(tree_array[0])
        return
    }
    current_node, index = get_leftmost_descendant(tree_array, index)
    while True {
        if current_node is None {
            raise ValueError
        }
        visit(current_node)
        if current_node.right_thread is None {
            return
        }
        if current_node.right_is_thread {
            index = current_node.right_thread
            current_node = tree_array[index]
        }
        else {
            tree_array, rc_index = right_child_index(tree_array, index)
            current_node, index = get_leftmost_descendant(tree_array, rc_index)
        }
    }
}
\end{minted}

\pagebreak
\noindent\textbf{Problem 3:}
\medskip

Let's sort using a method not discussed in class. Suppose you have $n$ data values in in array \texttt{A}.  Declare an array called \texttt{Count}.  Look at the value in \texttt{A[i]}. Count the number of items in A that are smaller than the value in \texttt{A[i]}.  Assign that result to \texttt{count[i]}. Declare an output array \texttt{Output}. Assign \texttt{Output[count[i]] = A[i]}. Think about what the size of \texttt{Output} needs to be. Is it $n$ or something else? Write a method to sort based on this strategy.

\medskip
\noindent\textbf{Answer:}
\medskip

\begin{minted}[linenos,breaklines,breakanywhere=true,fontsize={\fontsize{8pt}{9pt}\selectfont}]{Go}
function mystery_sort(A) {
    L = len(A)
    C = []
    C.extend([null]*L)
    output = []
    output.extend([null]*L)
    for i, n in enumerate(A) {
        sum = 0
        for j in A {
            if j < n {
                sum += 1
            }
            C[i] = sum
        }
    }
    for i in C {
        output[C[i]] = A[i]
    }
    for i, x in enumerate(output) {
        if x is None {
            output[i] = output[i-1]
        }
    }
    return output
}
\end{minted}

\medskip
\noindent\textbf{Problem 4:}
\medskip

Analyze the cost of the sort you wrote in the previous problem. What is the impact of random, ordered, or reverse ordered data?

\medskip
\noindent\textbf{Answer:}
\medskip


\medskip
\noindent\textbf{Problem 5:}
\medskip

How many comparisons are necessary to find the largest and smallest of a set of $n$ distinct elements? Do not assume the answer must involve sorting. It could but does not need to do so. Try to be as efficient as you can.

\medskip
\noindent\textbf{Answer:}
\medskip


\pagebreak

\section*{Source Code}
The full source code, including LaTeX, tooling, other artifacts, and git history is available on my Github repository for the class here:
\begin{center}
  \url{https://github.com/jllovet/jhu-en-605-202-su26-data-structures}
\end{center}

\printbibliography

\end{document}
